\documentclass[11pt]{article}

% ACL-formatted experimental report. ``preprint'' keeps the ACL two-column
% layout without review line numbers.
\usepackage[preprint]{acl}
\usepackage{times}
\usepackage{latexsym}
\usepackage[T1]{fontenc}
\usepackage[utf8]{inputenc}
\usepackage{microtype}
\usepackage{inconsolata}
\usepackage{booktabs}
\usepackage{array}
\usepackage{amsmath}
\usepackage{enumitem}
\usepackage{float}
\usepackage{graphicx}
\usepackage[section]{placeins}

% Keep experiment tables close to their explanation and pack float-only pages
% from the top if LaTeX must create one.
\makeatletter
\setlength{\@fptop}{0pt}
\setlength{\@fpsep}{12pt}
\setlength{\@fpbot}{0pt plus 1fil}
\makeatother
\renewcommand{\topfraction}{0.95}
\renewcommand{\bottomfraction}{0.95}
\renewcommand{\textfraction}{0.05}
\renewcommand{\floatpagefraction}{0.80}
\setlength{\textfloatsep}{12pt plus 2pt minus 2pt}

\newcommand{\gold}{\textsc{Golden Test}}
\newcommand{\Tzero}{\textsc{T0}}
\newcommand{\Tone}{\textsc{T1}}
\newcommand{\Ttwo}{\textsc{T2}}
\newcommand{\Rc}{\ensuremath{R_{c}}}
\newcommand{\Rd}{\ensuremath{R_{d}}}
\newcommand{\Judge}{\ensuremath{J}}
\newcommand{\best}[1]{\textbf{#1}}

\title{CV801 Experimental Report:\\
Detection, Grounded Agricultural Captioning, and Multi-Agent Revision}

\author{CV801 Research Project \\
Internal Experimental Report}

\begin{document}
\maketitle

\begin{abstract}
This report documents the completed CV801 experiment chain in execution order.
The testbed contains 8,097 agricultural images and 14,601 annotated objects,
including a frozen 1,000-image test set. Three detectors were fine-tuned and
evaluated under a unified COCO protocol. A 9B multimodal writer was then tested
zero-shot, after domain LoRA adaptation, and with predicted object proposals.
Finally, visual and domain reviewers, a text-only judge, one- and two-round
revision, an alternative reviewer, and an independent visual evaluator were
compared. Fine-tuned YOLO26x reached 72.9 COCO AP. Domain adaptation increased
caption token F1 from 5.8 to 49.1. Predicted proposals increased object-link F1
from 59.7 to 76.7, although the independent visual score changed only from
4.172 to 4.175 out of 6. One full review round produced the highest visual score
of 4.288; a second round regressed to 4.268 and consumed more GPU time. Human
rating and a direct GPT-6 Astra baseline were cancelled before finalization and
are not pending experiments.
\end{abstract}

\section{Scope and Completion Status}

The experimental goal is to generate one structured agricultural report per
image. The report must identify every visible target object, keep descriptions
linked to object IDs, provide a conservative assessment and confidence, and
return valid JSON. The completed chain addresses three questions:

\begin{enumerate}[leftmargin=*,nosep]
    \item Does domain adaptation improve structured descriptions over a
    zero-shot multimodal model?
    \item Do predicted object proposals improve object enumeration and linkage?
    \item Does role-separated Agent Review improve the final report enough to
    justify its compute cost?
\end{enumerate}

All retained automated runs, frozen-test evaluations, corrected-denominator
validation, unified COCO and operating-point detection evaluation, independent
visual evaluation, and per-image compute aggregation have completed. Human or
engineer scoring was removed from the final scope. The direct GPT-6 Astra
baseline was also cancelled. They are reported as scope decisions, not
incomplete work.

\section{Dataset and Common Protocol}

\subsection{Data Split}

The cleaned dataset contains 8,097 images and 14,601 accepted boxes from 11 crop
categories: bitter gourd, cauliflower, cherry tomato, chili, date palm,
eggplant, fig, lemon, luffa, mulberry, and tomato. Table~\ref{tab:data} shows the
fixed split used by every experiment.

\begin{table}[H]
\centering
\small
\begin{tabular}{lrr}
\toprule
Split & Images & Objects \\
\midrule
Train & 6,387 & 11,555 \\
Validation & 710 & 1,299 \\
\gold & 1,000 & 1,747 \\
\midrule
Total & 8,097 & 14,601 \\
\bottomrule
\end{tabular}
\caption{Frozen split shared by detection and report-generation experiments.}
\label{tab:data}
\end{table}

Splitting uses seed 20261104 and stratifies by crop and fruit/leaf context. It
keeps 7,862 source groups intact, with no group crossing train, validation, and
test. Ninety-three images whose boxes were all rejected during cleaning are
excluded instead of being treated as negative images. Images are converted to
EXIF-corrected RGB JPEGs for model input; the original files are unchanged.

Of the training images, 2,478 contain multiple accepted objects. Therefore, each
image appears once in the writer dataset and all targets are aggregated into one
ordered response. Objects are numbered from top to bottom and then left to
right. Supervised report targets are dataset-level visually reviewed GPT
captions, not independently written engineer references. This provenance also
applies to reference-based caption metrics.

\subsection{Frozen-Test and Leakage Controls}

The 710-image validation split selects models and protocols. The 1,000-image
\gold{} remains unavailable to training and selection. The formal test DAG runs
all inference without labels. References are attached only after predictions
and reviewer outputs are frozen. The writer, visual reviewer, domain reviewer,
text-only judge, and independent visual evaluator do not receive test reference
texts.

Every output is audited for expected sample count, unique sample IDs, JSON
validity, and schema compliance before aggregation. All formal training runs use
eight AMD Instinct MI210 GPUs with approximately 64~GiB memory per GPU, ROCm 7,
and PyTorch 2.8 containers.

\section{Experiment A: Object Detection}

\subsection{Models and Training Settings}

The detection experiment compares a closed-set detector with two grounding-
oriented models. YOLO26x is initialized from the Ultralytics checkpoint
\citep{ultralytics2026}. Grounding DINO Tiny uses language-guided detection
\citep{liu2023grounding}. LocateAnything-3B uses generative parallel box
decoding \citep{wang2026locate}. Table~\ref{tab:det-settings} records the formal
fine-tuning settings.

\begin{table}[H]
\centering
\scriptsize
\renewcommand{\arraystretch}{1.10}
\begin{tabular}{p{1.85cm}p{4.85cm}}
\toprule
Model & Formal setting \\
\midrule
YOLO26x & 100-epoch cap; early stop at epoch 88; patience 30; 960 px; global batch 64; AMP; cosine schedule; automatic optimizer; weight decay $5\times10^{-4}$; seed 20261104. \\
Grounding DINO Tiny & 12 epochs / 4,800 steps; batch 1 per GPU; accumulation 2; effective batch 16; LR $2\times10^{-5}$; BF16; cosine schedule; seed 20261104. \\
LocateAnything-3B & 3,000 steps; batch 1 per GPU; effective batch 8; LR $2\times10^{-5}$; warmup 100; BF16; SDPA; 4,096 tokens; packing buffer 32; LoRA rank 64; language and vision backbones frozen; connector trainable. \\
\bottomrule
\end{tabular}
\caption{Formal detector fine-tuning settings. All runs use eight MI210 GPUs.}
\label{tab:det-settings}
\end{table}

Before formal training, each model passes runtime, smoke-fit, and output-format
checks. YOLO's best framework-native validation checkpoint occurs at epoch 58,
with mAP50--95 of 0.624. This native validation number is not mixed with the
unified final COCO evaluation below.

\subsection{Unified COCO Results}

All detectors are evaluated on the same 1,000 \gold{} images using
\texttt{pycocotools}. We report COCO AP averaged over IoU 0.50--0.95 and AP50
\citep{lin2014coco}. YOLO uses native confidence scores, confidence threshold
0.001, and NMS IoU 0.7. Grounding DINO uses box and text thresholds 0.01 and NMS
IoU 0.6. LocateAnything does not expose a calibrated native confidence in the
evaluated interface, so every box receives score 0.2 and generation order
breaks ties.

\begin{table}[H]
\centering
\small
\begin{tabular}{lrrr}
\toprule
Detector & AP & AP50 & Boxes \\
\midrule
YOLO26x & \best{72.9} & \best{82.2} & 14,280 \\
Grounding DINO Tiny & 26.5 & 28.2 & 73,517 \\
LocateAnything-3B & 53.3 & 61.8 & 1,468 \\
\bottomrule
\end{tabular}
\caption{Unified COCO box metrics (\%) on the 1,000-image \gold.}
\label{tab:det-results}
\end{table}

The thresholded metrics requested in the experiment plan use class-aware,
one-to-one greedy matching at IoU $\geq0.5$. Thresholds are frozen without
using \gold{} labels: 0.404 for YOLO from its validation F1 curve, 0.250913 for
Grounding DINO by maximum validation micro-F1, and 0.2 for all LocateAnything
boxes because its interface assigns one constant score. These operating points
are separate from the low thresholds used to retain detections for COCO AP.

\begin{table}[H]
\centering
\small
\setlength{\tabcolsep}{3.2pt}
\resizebox{\columnwidth}{!}{%
\begin{tabular}{lrrrrr}
\toprule
Detector & Score & Prec. & Recall & F1 & TP-IoU \\
\midrule
YOLO26x & 0.404 & 76.1 & \best{80.5} & \best{78.2} & 0.943 \\
Grounding DINO Tiny & 0.251 & 22.5 & 59.5 & 32.7 & \best{0.971} \\
LocateAnything-3B & 0.200 & \best{85.3} & 71.7 & 77.9 & 0.950 \\
\bottomrule
\end{tabular}
}
\caption{Frozen operating-point micro metrics (\%) on \gold. TP-IoU is mean
IoU over class-correct true-positive matches and does not reward missed boxes.}
\label{tab:det-operating}
\end{table}

YOLO26x obtains the highest AP, exceeding LocateAnything by 19.6 points and
Grounding DINO by 46.4 points, and the highest operating-point F1 at 78.2. It is
therefore frozen as the proposal generator for the subsequent T2 writer.
LocateAnything reaches a similar F1 through higher precision and lower recall,
but its constant confidence limits interpretation of its ranked
precision--recall curve. Grounding DINO's matched boxes have high IoU, but its
many false positives reduce precision to 22.5\%.

\begin{table}[H]
\centering
\small
\setlength{\tabcolsep}{2.5pt}
\begin{tabular}{lrr@{\hspace{4pt}}lrr}
\toprule
Crop & AP & AP50 & Crop & AP & AP50 \\
\midrule
Bitter gourd & 74.3 & 85.6 & Fig & 74.1 & 80.8 \\
Cauliflower & 83.6 & 92.3 & Lemon & \best{87.4} & 93.6 \\
Cherry tomato & 84.3 & 92.3 & Luffa & 58.1 & 73.2 \\
Chili & 52.8 & 69.8 & Mulberry & 75.2 & 83.3 \\
Date palm & 43.9 & 48.1 & Tomato & 83.0 & 91.6 \\
Eggplant & 84.9 & 93.9 & & & \\
\bottomrule
\end{tabular}
\caption{YOLO26x per-category COCO metrics (\%) on \gold.}
\label{tab:per-category}
\end{table}

Per-category AP ranges from 43.9 for date palm to 87.4 for lemon. Chili and
luffa also fall below 60 AP. The detector is strongest overall but does not
provide uniform crop coverage; these lower-performing classes remain a boundary
for downstream proposal conditioning.

\FloatBarrier

\section{Experiment B: Structured VLM Writer}

\subsection{Compared Writer Configurations}

The writer returns image-level \texttt{crop} and \texttt{plant\_part} fields and
an ordered object list. Each object contains an ID, caption, assessment, image
quality, confidence, and action. Coordinates are never copied into the output.
The four writer configurations are:

\begin{itemize}[leftmargin=*,nosep]
    \item \Tzero: zero-shot Qwen3.5-9B \citep{qwen2026qwen35}, whole image, no boxes;
    \item \Tone: Qwen3.5-9B with domain LoRA, whole image, no boxes;
    \item \Ttwo: the same domain training and targets as \Tone, with frozen
    YOLO predictions added to the prompt; and
    \item Qwen3.5-27B: an unadapted, zero-shot scale baseline, whole image, no boxes.
\end{itemize}

For \Ttwo, proposals are normalized to $[0,1000]$ \texttt{xyxy} coordinates,
filtered at score 0.25, and capped at 20 per image. The prompt explicitly says
that they may be missing or incorrect and must be treated as localization hints,
not labels. T1 and T2 assistant targets are checked for exact identity, so the
comparison isolates predicted-proposal conditioning.

\begin{table}[H]
\centering
\scriptsize
\renewcommand{\arraystretch}{1.08}
\setlength{\tabcolsep}{2.5pt}
\begin{tabular}{p{2.25cm}p{2.15cm}p{2.15cm}}
\toprule
Setting & \Tone & \Ttwo \\
\midrule
Training examples & 6,387 / 710 & Same \\
Input & Whole image & Image + proposals \\
Epochs / steps & 3 / 1,197 & 3 / 1,197 \\
Batch/GPU; accum.; GPUs & 1; 2; 8 & 1; 2; 8 \\
Effective batch & 16 & 16 \\
Learning rate & $10^{-4}$ & $10^{-4}$ \\
Schedule / warmup & Cosine / 5\% & Same \\
Weight decay & 0.01 & 0.01 \\
LoRA & $r=64$, $\alpha=128$, dropout 0.05 & Same \\
Trainable modules & Language all-linear LoRA & Same \\
Frozen modules & Vision tower + aligner & Same \\
Precision / attention & BF16 / SDPA & Same \\
Training length & 4,096 tokens & 6,144 tokens \\
Image-token maximum & 768 & 768 \\
Inference & temp. 0; 4,096 new tokens & Same \\
\bottomrule
\end{tabular}
\caption{Formal T1 and T2 training settings. T0 and the 27B baseline receive no
fine-tuning. LoRA follows \citet{hu2022lora}.}
\label{tab:writer-settings}
\end{table}

The 6,144-token T2 limit covers the proposal-plus-target long tail. The longest
audited targets contain 35 objects in training and 22 in validation. T1 reaches
validation loss 0.483 and token accuracy 0.841; T2 reaches 0.491 and 0.841.

\subsection{Evaluation Metrics}

Structured metrics include valid JSON, exact object count, and object-ID link
precision, recall, and F1. Text metrics compare only unique object-ID matches:
caption token F1, unsmoothed single-reference corpus BLEU-4
\citep{papineni2002bleu}, and mean object-level ROUGE-L F1
\citep{lin2004rouge}. These metrics measure agreement with the reviewed GPT
targets, not independent human correctness.

Field-level metrics compare crop, plant part, image quality, stage, disease,
severity, action, and confidence with the same reviewed GPT targets. Confidence
is reported as mean absolute error (MAE; lower is better); categorical fields
and action are exact-match or token-F1 measures. Missing fields are reported as
unavailable rather than scored as zero.

A separate Pixtral-12B evaluator \citep{agrawal2024pixtral} sees the image and
candidate report without reference text. It assigns integer scores from 0 to 2
for visual evidence ($S_e$), conclusion support ($S_s$), and confidence
appropriateness ($S_c$):

\begin{equation}
S_{\mathrm{visual}}=S_e+S_s+S_c, \qquad S_{\mathrm{visual}}\in[0,6].
\end{equation}

This independent score is model-based and has not been calibrated against human
agreement.

\subsection{Writer Results}

\begin{table}[H]
\centering
\scriptsize
\setlength{\tabcolsep}{2.2pt}
\resizebox{\columnwidth}{!}{%
\begin{tabular}{lrrrrrrrr}
\toprule
System & JSON & Count & ID-P & ID-R & ID-F1 & Cap.-F1 & B-4 & R-L \\
\midrule
\Tzero & 97.6 & 57.2 & 43.1 & 56.3 & 48.9 & 5.8 & 0.0 & 5.7 \\
\Tone & 88.8 & 52.9 & 56.8 & 62.9 & 59.7 & 49.1 & 14.0 & 39.0 \\
\Ttwo & \best{99.5} & \best{64.1} & 67.2 & \best{89.3} & \best{76.7} & \best{51.3} & \best{15.5} & \best{42.2} \\
27B zero-shot & 98.4 & 60.4 & \best{99.2} & 56.5 & 72.0 & 22.2 & 0.4 & 17.4 \\
\bottomrule
\end{tabular}
}
\caption{Writer structure and text results (\%) on \gold. B-4 is BLEU-4 and R-L
is ROUGE-L.}
\label{tab:writer-structured}
\end{table}

Domain adaptation is the main language intervention. Relative to T0, T1 raises
caption F1 from 5.8 to 49.1 and ROUGE-L from 5.7 to 39.0. The larger 27B
zero-shot model reaches only 22.2 caption F1, so parameter scale alone does not
replace in-domain adaptation in this comparison.

Predicted boxes primarily improve structure. Relative to T1, T2 raises exact
object count by 11.2 points, object-link F1 by 17.0 points, and caption F1 by 2.2
points. JSON validity also increases from 88.8\% to 99.5\%.

\begin{table}[H]
\centering
\scriptsize
\resizebox{\columnwidth}{!}{%
\begin{tabular}{lrrrrrrrr}
\toprule
System & Crop & Part & ImgQ & Stage & Disease & Severity & Action & Conf.-MAE \\
\midrule
\Tzero & 97.6 & 97.6 & \best{51.9} & -- & -- & -- & 0.6 & -- \\
\Tone & 88.8 & 88.8 & 51.4 & \best{73.2} & 40.4 & 39.3 & \best{42.0} & \best{0.132} \\
\Ttwo & \best{99.4} & \best{99.5} & 49.4 & 60.8 & \best{42.2} & \best{41.0} & 37.0 & 0.145 \\
27B zero-shot & 63.6 & 93.7 & 17.6 & -- & -- & -- & 3.8 & 0.203 \\
\bottomrule
\end{tabular}
}
\caption{Writer field-level agreement with reviewed GPT targets (\%; confidence
is MAE, lower is better). Dashes mean that the output omitted the field.}
\label{tab:writer-fields}
\end{table}

The field audit prevents the structural gain from being read as a uniform
semantic gain. T2 is better than T1 for crop, plant part, disease, and severity,
but T1 is better for stage, action, and confidence MAE. Proposal conditioning
therefore improves enumeration while trading off some report fields.

\begin{table}[H]
\centering
\small
\begin{tabular}{lrrrr}
\toprule
System & Evidence & Support & Conf. & Total \\
\midrule
\Tzero & 1.285 & 0.560 & 0.805 & 2.650 \\
\Tone & \best{1.623} & 1.275 & 1.274 & 4.172 \\
\Ttwo & 1.512 & \best{1.353} & \best{1.310} & \best{4.175} \\
27B zero-shot & 1.116 & 0.730 & 0.982 & 2.828 \\
\bottomrule
\end{tabular}
\caption{Independent model-based visual score for writer configurations. Each
component is in $[0,2]$ and Total is in $[0,6]$.}
\label{tab:writer-visual}
\end{table}

The independent visual result narrows the interpretation. T2 improves support
and confidence appropriateness over T1, but its visual-evidence component drops
from 1.623 to 1.512. The total is almost unchanged at 4.172 versus 4.175.
Predicted proposals therefore act as strong enumeration and linkage anchors;
this experiment does not show a general improvement in visual faithfulness.

\begin{table}[H]
\centering
\scriptsize
\setlength{\tabcolsep}{3.0pt}
\resizebox{\columnwidth}{!}{%
\begin{tabular}{lrrrrrr}
\toprule
System & Supported & Unsupported & Unsup.\% & Correct & Incorrect & Uncertain \\
\midrule
\Tzero & 1,349 & 2,194 & 61.9 & 637 & 217 & 151 \\
\Tone & 2,845 & 1,232 & \best{30.2} & 1,029 & 54 & 37 \\
\Ttwo & \best{3,266} & 1,621 & 33.2 & \best{1,142} & 104 & 137 \\
27B zero-shot & 1,528 & 1,994 & 56.6 & 506 & 341 & 161 \\
\bottomrule
\end{tabular}
}
\caption{Independent evaluator claim and object-link audit. Unsupported share
uses supported plus unsupported enumerated claims as its denominator.}
\label{tab:writer-claim-audit}
\end{table}

The model-based claim audit gives T1 the lowest unsupported-claim share at
30.2\%. T2 produces more supported claims than T1 (3,266 versus 2,845), but also
more unsupported claims (1,621 versus 1,232), raising the share to 33.2\%.
Thus, the stronger T2 object linkage does not establish lower claim-level
hallucination.

The count-based extra-object proxy falls from 56.9\% for T0 to 43.1\% for T1
and 32.8\% for T2. It is not a visual hallucination rate: it cannot detect an
unsupported disease, severity, action, or attribute, and cannot distinguish a
model error from a missing reference object.

\FloatBarrier

\section{Experiment C: Agent Review and Revision}

\subsection{Review Roles and Configurations}

T2 is frozen as the writer. Review follows the iterative-feedback motivation of
Self-Refine \citep{madaan2023selfrefine}, but separates roles and evidence
access:

\begin{itemize}[leftmargin=*,nosep]
    \item \Rc, InternVL3.5-8B \citep{wang2025internvl}, sees the image and draft
    and checks visible evidence, target coverage, and object linkage;
    \item \Rd, InternVL3.5-38B, sees the image and draft and checks agricultural
    terminology, maturity, disease, severity, confidence, and action; and
    \item \Judge, unadapted Qwen3.5-9B, sees only the draft and reviews. It can
    choose \texttt{accept} or \texttt{revise} but cannot add visual facts.
\end{itemize}

For the original matrix, the planned Gemma 3 visual reviewer was unavailable
behind an external model-access gate. InternVL3.5-8B was therefore the primary
\Rc{}, and Pixtral-12B provided the cross-family swap. Access became available
for the later additive reviewer-capacity experiment, which evaluates Gemma 3
reviewers without replacing the original results. \Rd{} uses its frozen model
knowledge but no external domain-reference corpus; all validation and test
reference texts remain hidden.

The compared configurations are Rd-only, Rc-only, Rc+Rd, the full
Rc+Rd+Judge pipeline after one revision, and the full pipeline after a second
revision and recheck. A reviewer-swap experiment replaces Rc with Pixtral-12B.
Revisions always use the frozen T2 LoRA writer. References remain hidden from
all agents.

\subsection{Structured Results After Revision}

\begin{table}[H]
\centering
\scriptsize
\resizebox{\columnwidth}{!}{%
\begin{tabular}{lrrrrrrrr}
\toprule
Configuration & JSON & Count & ID-P & ID-R & ID-F1 & Cap.-F1 & Extra \\
\midrule
Frozen \Ttwo & \best{99.5} & \best{64.1} & 67.2 & \best{89.3} & \best{76.7} & \best{51.3} & 32.8 \\
\Rd-only & 95.6 & 59.2 & \best{99.2} & 56.2 & 71.7 & 50.9 & \best{0.5} \\
\Rc-only & 93.8 & 56.7 & \best{99.2} & 54.7 & 70.5 & \best{51.3} & 0.7 \\
\Rc+\Rd & 95.6 & 59.4 & 84.3 & 56.5 & 67.6 & 50.1 & 15.5 \\
Full round 1 & 91.6 & 58.3 & 97.9 & 55.5 & 70.9 & \best{51.3} & 1.9 \\
Full round 2 & 97.8 & 59.6 & 89.1 & 59.6 & 71.4 & 50.4 & 10.9 \\
Pixtral as \Rc & 91.1 & 57.6 & 97.6 & 56.2 & 71.3 & 51.1 & 2.2 \\
\bottomrule
\end{tabular}
}
\caption{Structured review results (\%). Extra is the extra-object proxy, not a
visual hallucination metric.}
\label{tab:review-structured}
\end{table}

Every reviewed configuration lowers exact count and object-link F1 relative to
the frozen T2 writer. The very low extra-object proxy in several review variants
does not imply that they are visually correct. Together with the reduced count
metrics, it indicates that reviewers often make the output more conservative by
deleting or collapsing difficult objects.

\begin{table}[H]
\centering
\scriptsize
\resizebox{\columnwidth}{!}{%
\begin{tabular}{lrrrrrrrr}
\toprule
Configuration & Crop & Part & ImgQ & Stage & Disease & Severity & Action & Conf.-MAE \\
\midrule
Frozen \Ttwo & \best{99.4} & \best{99.5} & 49.4 & 60.8 & 42.2 & 41.0 & 37.0 & 0.145 \\
\Rd-only & 76.9 & 95.2 & 52.6 & \best{62.6} & 41.0 & 38.4 & 40.3 & 0.135 \\
\Rc-only & 72.2 & 93.5 & 52.9 & 58.7 & 44.1 & 42.6 & \best{41.1} & 0.133 \\
\Rc+\Rd & 66.5 & 95.1 & 50.9 & 58.5 & 41.8 & 39.8 & 40.3 & 0.132 \\
Full round 1 & 79.2 & 90.0 & 56.3 & 61.1 & 44.2 & \best{43.0} & 41.0 & 0.135 \\
Full round 2 & \best{85.7} & \best{95.5} & 52.8 & 60.2 & 43.3 & 41.4 & 39.4 & \best{0.131} \\
Pixtral as \Rc & 78.9 & 89.6 & \best{56.7} & 61.2 & \best{44.3} & 42.7 & 40.1 & 0.135 \\
\bottomrule
\end{tabular}
}
\caption{Field-level review results against reviewed GPT targets (\%; confidence
is MAE, lower is better).}
\label{tab:review-fields}
\end{table}

Review improves selected fields but not one configuration dominates. Full
round 1 is strongest on severity, Pixtral on image quality and disease, and
round 2 on confidence MAE. These gains occur alongside lower object count and
linkage recall than frozen T2.

\subsection{Independent Visual Results}

\begin{table}[H]
\centering
\small
\setlength{\tabcolsep}{3.7pt}
\begin{tabular}{lrrrr}
\toprule
Configuration & Evidence & Support & Conf. & Total \\
\midrule
Frozen \Ttwo & 1.512 & 1.353 & 1.310 & 4.175 \\
\Rd-only & 1.508 & 1.376 & 1.334 & 4.218 \\
\Rc-only & 1.531 & 1.368 & 1.343 & 4.242 \\
\Rc+\Rd & 1.488 & 1.373 & 1.278 & 4.139 \\
Full round 1 & \best{1.581} & 1.352 & \best{1.355} & \best{4.288} \\
Full round 2 & 1.559 & \best{1.387} & 1.322 & 4.268 \\
\bottomrule
\end{tabular}
\caption{Independent Pixtral-12B scores for review configurations. Pixtral as
Rc is omitted to prevent evaluator self-scoring.}
\label{tab:review-visual}
\end{table}

Rc-only and Rd-only improve the total from 4.175 to 4.242 and 4.218,
respectively. Combining the two reviews without the full judge loop falls to
4.139. The full first round is best at 4.288, a gain of 0.113 over T2. The
second round increases conclusion support but lowers evidence and confidence
components, producing a slightly lower total of 4.268. The additional revision
therefore does not improve the aggregate result.

\begin{table}[H]
\centering
\scriptsize
\setlength{\tabcolsep}{2.7pt}
\resizebox{\columnwidth}{!}{%
\begin{tabular}{lrrrrrr}
\toprule
Configuration & Supported & Unsupported & Unsup.\% & Correct & Incorrect & Uncertain \\
\midrule
Frozen \Ttwo & 3,266 & 1,621 & \best{33.2} & \best{1,142} & 104 & 137 \\
\Rd-only & 2,324 & 1,262 & 35.2 & 928 & 51 & 46 \\
\Rc-only & 2,330 & 1,257 & 35.0 & 927 & 49 & 43 \\
\Rc+\Rd & 2,221 & 1,254 & 36.1 & 927 & \best{43} & 45 \\
Full round 1 & \best{2,489} & 1,305 & 34.4 & 947 & 50 & 58 \\
Full round 2 & 2,370 & 1,278 & 35.0 & 940 & 58 & \best{27} \\
\bottomrule
\end{tabular}
}
\caption{Independent claim and object-link audit for review configurations.
Pixtral as \Rc{} is excluded to avoid evaluator self-scoring.}
\label{tab:review-claim-audit}
\end{table}

The higher 0--6 score after full round 1 is not accompanied by a lower
unsupported-claim share: 34.4\% versus 33.2\% for frozen T2. The revision also
reduces the number of independently audited links. The measured gain is
therefore narrow---higher evidence and confidence rubric components---rather
than a general reduction in model-flagged unsupported claims.

\subsection{Decision Behavior}

\begin{table}[H]
\centering
\small
\begin{tabular}{lrr}
\toprule
Stage & Accept/pass & Revise \\
\midrule
\Rc review & 24 & 976 \\
\Rd review & 235 & 765 \\
Text judge & 27 & 973 \\
Recheck round 1 & 29 & 971 \\
Recheck round 2 & 31 & 969 \\
Pixtral review & 0 & 1,000 \\
Pixtral judge & 2 & 998 \\
Pixtral recheck & 3 & 997 \\
\bottomrule
\end{tabular}
\caption{Review decisions over 1,000 \gold{} examples. Recheck ``Revise'' means
\texttt{revise\_again}.}
\label{tab:decisions}
\end{table}

The decision rates are much stricter than the measured gains. Rc requests
revision for 97.6\% of drafts, and the judge requests revision for 97.3\%. Even
after two rewrites, the judge still requests another revision for 96.9\%.
Pixtral is stricter: it requests revision for every initial draft and its judge
accepts only two. These distributions show that the current judge is not a
calibrated stopping rule.

\FloatBarrier

% Include the additive Writer-scale section only when the complete paired-grid
% tables exist. Earlier report sections and results remain unchanged.
\IfFileExists{writer_reviewer_grid_tables.tex}{%
  % Additive CV801 Writer-scale report section.
% Writer-only evidence was frozen from writer_results.current.csv. The
% reviewer-grid tables were added after the complete nine-cell paired audit and
% 10,000-resample bootstrap analysis passed on 2026-11-17 UAE time.

\section{Small-Writer Scaling and Reviewer-Capacity Expansion}

\subsection{Controlled Writer-size comparison}

We extended the Writer experiment with Qwen3.5-2B and Qwen3.5-4B while
retaining Qwen3.5-9B as the adapted anchor. For each controlled size, T0 is the
unadapted whole-image model, T1 uses the same three-epoch domain LoRA protocol,
and T2 adds the same frozen YOLO26x localization hints. Data, prompts, target
responses, decoding temperature, the 4,096-token Golden Test generation cap,
and evaluation code were held fixed. The existing 27B zero-shot result remains
a scale reference rather than a fourth point on the adapted scaling curve.

\begin{table}[H]
\centering
\scriptsize
\setlength{\tabcolsep}{2.0pt}
\resizebox{\columnwidth}{!}{%
\begin{tabular}{llrrrrrrrr}
\toprule
Writer & Variant & JSON & Schema & Count & ID-F1 & Cap.-F1 & Visual & Unsup. \\
\midrule
2B & T0 & 82.4 & 0.0 & 45.4 & 35.5 & 20.7 & 3.429 & 44.6 \\
2B & T1 & 0.0 & 0.0 & 0.0 & -- & -- & 3.588 & 32.9 \\
2B & T2 & 0.1 & 0.1 & 0.1 & 0.0 & 31.3 & 3.505 & 41.0 \\
4B & T0 & 69.4 & 0.0 & 4.3 & 27.1 & 16.7 & 3.078 & 44.6 \\
4B & T1 & 0.0 & 0.0 & 0.0 & -- & -- & 3.248 & 56.3 \\
4B & T2 & 5.2 & 5.2 & 4.5 & 5.8 & 45.5 & 4.631 & 33.7 \\
9B & T0 & 97.6 & 0.0 & 57.2 & 48.9 & 5.8 & 2.650 & 61.9 \\
9B & T1 & 88.8 & 88.8 & 52.9 & 59.7 & 49.1 & 4.172 & 30.2 \\
9B & T2 & 99.5 & 99.5 & 64.1 & 76.7 & 51.3 & 4.175 & 33.2 \\
\bottomrule
\end{tabular}%
}
\caption{Controlled Writer-size results on the 1,000-image Golden Test. JSON,
Schema, Count, ID-F1, Cap.-F1, and Unsup. are percentages. Visual is the
independent model-based score in $[0,6]$. Unsup. is the unsupported-claim share,
not a human hallucination rate. A dash indicates that the frozen structured
evaluator could not define the metric from schema-invalid outputs.}
\label{tab:small-writer-results}
\end{table}

The adapted small Writers did not preserve the structured-output reliability
of the 9B Writer under the frozen protocol. Both 2B-T1 and 4B-T1 produced no
schema-valid reports. Proposal conditioning recovered only 0.1\% schema
validity for 2B-T2 and 5.2\% for 4B-T2, compared with 99.5\% for 9B-T2.
Object-ID link F1 showed the same capacity boundary: 0.0\% for 2B-T2 and 5.8\%
for 4B-T2 versus 76.7\% for 9B-T2. The controlled comparison therefore does
not support replacing the 9B Writer with either smaller adapted model when
machine-readable output is required.

The independent visual score alone gives a different and potentially
misleading ranking. In particular, 4B-T2 reached 4.631, above the 9B-T2 score
of 4.175, even though only 5.2\% of its reports satisfied the output schema.
This divergence shows that the visual rubric can reward locally plausible
content while missing an end-to-end interface failure. We therefore treat
schema validity as a non-negotiable deployment constraint and interpret the
visual score only jointly with structured metrics. The result also explains
why invalid samples remain in the evaluation denominator rather than being
dropped.

These findings are specific to the fixed optimization and decoding protocol.
Each adapted size uses one training seed, so the experiment identifies a
failure boundary for the tested checkpoints rather than proving that 2B or 4B
models cannot be stabilized under any size-specific tuning regime.

\begin{table}[H]
\centering
\small
\setlength{\tabcolsep}{4pt}
\begin{tabular}{llrrr}
\toprule
Writer & Variant & GPUs & GPU-h & s/image \\
\midrule
2B & T0 & 4 & 9.896 & 8.906 \\
2B & T1 & 4 & 25.836 & 23.252 \\
2B & T2 & 4 & 5.170 & 4.653 \\
4B & T0 & 4 & 29.606 & 26.645 \\
4B & T1 & 4 & 20.423 & 18.381 \\
4B & T2 & 8 & 84.989 & 38.245 \\
9B & T0 & 8 & 8.518 & 3.833 \\
9B & T1 & 8 & 40.400 & 18.180 \\
9B & T2 & 8 & 8.993 & 4.047 \\
\bottomrule
\end{tabular}
\caption{Measured Writer inference allocation for the controlled scale
extension. GPU-hours are elapsed allocation time multiplied by requested GPU
count; seconds per image use the 1,000-image Golden Test. These rows exclude
shared frozen-detector proposal generation and downstream review/evaluation.
The 9B jobs completed their predictions but exited during the superseded
post-inference audit; their complete outputs passed the corrected audit.}
\label{tab:small-writer-efficiency}
\end{table}

\subsection{Writer by reviewer-capacity interaction}

The second expansion crossed the T2 Writers with three reviewer tiers:
R-low (Gemma~3 4B \Rc{} and InternVL3.5-8B \Rd{}), R-mid (Gemma~3 12B
\Rc{} and InternVL3.5-14B \Rd{}), and R-high (Gemma~3 27B \Rc{} and
InternVL3.5-38B \Rd{}). The text-only Judge remained Qwen3.5-9B, and revision
used the same Writer size and T2 adapter as the initial draft. Thus, stronger
reviewers could not be credited with a hidden switch to the 9B revision model.
All configurations used the same 1,000 Golden Test images and the same
independent Pixtral-12B evaluator.

The primary analysis tests whether R-high closes the quality gap between a
small Writer and the 9B Writer relative to R-low. It reports the independent
visual total, object-ID link F1, schema validity, and unsupported-claim share.
For each endpoint, the two prespecified gap-closure contrasts are

\begin{align}
\Delta_{2B} &= Y_{2B,\mathrm{high}}-Y_{9B,\mathrm{high}} \notag\\
             &\quad -(Y_{2B,\mathrm{low}}-Y_{9B,\mathrm{low}}),\\
\Delta_{4B} &= Y_{4B,\mathrm{high}}-Y_{9B,\mathrm{high}} \notag\\
             &\quad -(Y_{4B,\mathrm{low}}-Y_{9B,\mathrm{low}}).
\end{align}

Uncertainty is estimated with 10,000 paired image-cluster bootstrap resamples
using seed 20261104. Object-level sufficient statistics are recomputed within
each resample so that multiple objects from the same image are not treated as
independent observations. A cell is eligible only after it contains exactly
1,000 unique and identically ordered Golden Test IDs. Invalid model outputs
remain in the denominator and receive the frozen endpoint-specific handling.

\IfFileExists{writer_reviewer_grid_tables.tex}{%
  % Auto-generated from paired_analysis.json; do not edit by hand.
% Frozen protocol: 1,000 paired images, 10,000 bootstrap resamples, seed 20261104.
\begin{table}[H]
\centering
\small
\setlength{\tabcolsep}{5pt}
\resizebox{\columnwidth}{!}{%
\begin{tabular}{llrrr}
\toprule
Endpoint & T2 Writer & R-low & R-mid & R-high \\
\midrule
Independent visual total & 2B & 3.568 & 3.597 & 3.548 \\
 & 4B & 3.196 & 3.176 & 3.223 \\
 & 9B & 4.240 & 4.265 & 4.212 \\
\addlinespace
Object-ID link F1 & 2B & 0.2\% & 0.6\% & 1.4\% \\
 & 4B & -- & -- & -- \\
 & 9B & 72.5\% & 73.5\% & 71.5\% \\
\addlinespace
Schema-valid rate & 2B & 6.6\% & 10.2\% & 14.2\% \\
 & 4B & 0.0\% & 0.0\% & 0.0\% \\
 & 9B & 92.8\% & 94.3\% & 91.0\% \\
\addlinespace
Unsupported-claim share & 2B & 46.3\% & 45.5\% & 45.1\% \\
 & 4B & 48.8\% & 47.8\% & 47.1\% \\
 & 9B & 34.9\% & 34.0\% & 34.6\% \\
\bottomrule
\end{tabular}
}%
\caption{Primary endpoints for the controlled Writer $\times$ reviewer grid. All nine cells contain the same 1,000 Golden Test images. Rates and F1 are percentages; independent visual total is reported on its 0--6 scale.}
\label{tab:writer-reviewer-grid}
\end{table}

\begin{table}[H]
\centering
\scriptsize
\setlength{\tabcolsep}{3.5pt}
\resizebox{\columnwidth}{!}{%
\begin{tabular}{llrr}
\toprule
Endpoint & Gap closure & Estimate & Paired-bootstrap 95\% CI \\
\midrule
Independent visual total & 2B vs. 9B & +0.008 & [-0.100, +0.113] \\
 & 4B vs. 9B & +0.055 & [-0.086, +0.193] \\
\addlinespace
Object-ID link F1 & 2B vs. 9B & +2.2 pp & [+0.6 pp, +3.7 pp] \\
 & 4B vs. 9B & -- & -- \\
\addlinespace
Schema-valid rate & 2B vs. 9B & +9.4 pp & [+6.3 pp, +12.6 pp] \\
 & 4B vs. 9B & +1.8 pp & [-0.1 pp, +3.7 pp] \\
\addlinespace
Unsupported-claim share & 2B vs. 9B & -0.8 pp & [-3.3 pp, +1.5 pp] \\
 & 4B vs. 9B & -1.3 pp & [-3.5 pp, +0.9 pp] \\
\bottomrule
\end{tabular}
}%
\caption{Prespecified high-minus-low reviewer gap-closure contrasts. Positive values use the raw endpoint direction. For unsupported-claim share, which is lower-is-better, a negative estimate is favorable. Intervals use 10,000 paired image-cluster bootstrap resamples. A dash means that the frozen metric is undefined because a required denominator is zero.}
\label{tab:writer-reviewer-interaction}
\end{table}
%
}{%
  \begin{quote}
  \textit{The four-endpoint Writer $\times$ reviewer table and the prespecified
  paired-bootstrap interaction table are inserted automatically only after all
  nine cells pass the identical 1,000-image audit. No partial-cell values are
  treated as reportable results.}
  \end{quote}
}

Reviewer capacity did not produce a clear interaction in the independent
visual total. The 2B and 4B gap-closure estimates were $+0.008$ (95\% CI
$[-0.100,0.113]$) and $+0.055$ ($[-0.086,0.193]$), respectively. The absolute
9B scores remained higher in every tier (4.212--4.265) than the 2B
(3.548--3.597) or 4B (3.176--3.223) scores.

The structured endpoints expose a narrower 2B benefit but no small-Writer
replacement. From R-low to R-high, 2B schema validity increased from 6.6\% to
14.2\%, whereas the 9B value changed from 92.8\% to 91.0\%; the direct
gap-closure estimate was $+9.4$ percentage points ($[+6.3,+12.6]$). The 2B
object-link interaction was $+2.2$ points ($[+0.6,+3.7]$), but its absolute
R-high F1 was still only 1.4\% versus 71.5\% for 9B. All three 4B revised
outputs had 0\% schema validity and no predicted objects under the frozen
parser, so their object-link F1 and the corresponding interaction are
undefined rather than zero. Unsupported-claim interactions were small and
their intervals crossed zero for both 2B ($-0.8$ points,
$[-3.3,+1.5]$) and 4B ($-1.3$ points, $[-3.5,+0.9]$). Thus the controlled
grid supports a modest reviewer-associated recovery for the 2B structured
interface, but stronger reviewers did not overcome the large absolute Writer
capacity gap under this protocol.

\subsection{Measured efficiency and interpretation boundary}

The final cost analysis uses allocated GPU-hours, execution span excluding
queue wait, seconds per image, and peak allocated memory when available. Shared
stages follow the charging rule used in the original report. Parameter count
is configuration metadata rather than a proxy for measured cost. The final
Pareto analysis will retain only observed quality--compute trade-offs and will
not label a small Writer efficient merely because its initial draft is cheaper.

The expansion has two interpretation boundaries. First, the reviewer tiers
change \Rc{} and \Rd{} jointly; the separate one-factor reviewer sweeps are
needed before attributing a tier effect to either role. Second, an improvement
in the model-based visual total cannot compensate for near-zero schema validity.
The complete conclusion must therefore combine interface validity, object
linkage, independent visual support, unsupported-claim share, and measured
compute rather than selecting the most favorable aggregate score.
%
  \FloatBarrier
}{}

\section{Efficiency and Runtime}

GPU-hours sum allocated GPU time for each logical system. Execution span is the
difference between its first start and final end, excluding scheduler wait.
Shared stages are charged to every logical system that requires them. The table
therefore documents the actual experiment chains, not perfectly cached marginal
serving cost. Because the internal MI210 allocation has no project-specific
monetary rate card, cost is reported as GPU-hours rather than an invented dollar
amount.

\begin{table}[H]
\centering
\scriptsize
\resizebox{\columnwidth}{!}{%
\begin{tabular}{lrrrrr}
\toprule
Logical system & GPU-h & GPU-h/img & s/img & Span (min) & Visual \\
\midrule
YOLO test inference & 0.149 & 0.000149 & 0.067 & 1.1 & -- \\
Grounding DINO test & 0.836 & 0.000836 & 0.376 & 6.3 & -- \\
LocateAnything test & 0.876 & 0.000876 & 0.394 & 6.6 & -- \\
\Tzero writer & 8.518 & 0.008518 & 3.833 & 63.9 & 2.650 \\
\Tone chain & 40.400 & 0.040400 & 18.180 & 303.0 & 4.172 \\
\Ttwo test chain & 9.143 & 0.009143 & 4.511 & 75.2 & 4.175 \\
27B zero-shot & 16.853 & 0.016853 & 7.584 & 126.4 & 2.828 \\
\Rc-only & 16.065 & 0.016065 & 53.441 & 890.7 & 4.242 \\
\Rd-only & 33.757 & 0.033757 & 72.250 & 1,204.2 & 4.218 \\
\Rc+\Rd & 37.719 & 0.037719 & 59.862 & 997.7 & 4.139 \\
Full round 1 & 44.999 & 0.044999 & 119.715 & 1,995.3 & \best{4.288} \\
Full round 2 & 51.523 & 0.051523 & 144.308 & 2,405.1 & 4.268 \\
Pixtral reviewer chain & 46.079 & 0.046079 & 119.778 & 1,996.3 & --$^{\dagger}$ \\
\bottomrule
\end{tabular}
}
\caption{Complete allocation accounting. $^{\dagger}$Pixtral reviewer output is
not scored by the same Pixtral evaluator. Detector rows are test inference, not
training cost.}
\label{tab:efficiency}
\end{table}

The full first review round consumes 45.0 GPU-hours, about 4.9 times the T2 test
chain, for a 0.113-point gain on the 0--6 visual score. The second round adds 6.5
GPU-hours and lowers the score by 0.020. For this dataset, frozen T2 is the
efficient choice when object structure matters; full round 1 is the higher-cost
choice when the small model-based visual-score gain is required.

\section{Consolidated Findings}

\paragraph{Detection.}
Fine-tuned YOLO26x is the selected locator at 72.9 COCO AP. Performance remains
crop-dependent, with date palm, chili, and luffa forming the lower tail. At its
validation-frozen operating point, it reaches 76.1\% precision, 80.5\% recall,
and 78.2\% F1.

\paragraph{Domain adaptation.}
T1 supplies the largest language improvement. Caption F1 rises from 5.8 to
49.1, while the 27B zero-shot baseline reaches only 22.2. In this setting,
in-domain supervision is more useful than increasing zero-shot model scale.

\paragraph{Predicted proposals.}
T2 improves exact object count and object-link F1 substantially, but its
independent visual total is almost tied with T1. The proposals improve
enumeration and association, not every claim's visual support.

\paragraph{Agent Review.}
One full round produces the best independent total of 4.288, but reduces count
and linkage metrics and has high cost. A second round does not help. Review
routing should therefore be selective rather than mandatory for every image.

\section{Interpretation Boundaries}

The results come from one agricultural collection and 11 crop categories.
Geographic, seasonal, sensor, and farm-level generalization are not measured.
Each configuration is represented by one frozen checkpoint on one fixed split;
the tables are descriptive and do not estimate multi-seed variance or confidence
intervals.
Reference captions are reviewed GPT outputs rather than independent agronomist
ground truth. The independent 0--6 score uses one visual model and has no human
calibration. Human scoring was cancelled, so the report makes no inter-rater or
expert-agreement claim. LocateAnything's constant score limits its AP ranking
interpretation. Operating-point P/R/F1 values use different validation-frozen
score policies and should not replace threshold-free AP. Unsupported-claim
shares are Pixtral enumerations rather than human hallucination labels. Cost
accounting follows actual scheduled chains and shared-stage charging rather
than optimized serving.

The system is decision support, not an autonomous treatment system. Generated
disease, severity, confidence, and action fields may be unsupported even when
object count is correct. Agronomist verification is required before acting on
intervention recommendations.

\section{Conclusion}

The completed experiments separate three effects. Domain LoRA is responsible
for most of the language-quality gain. Frozen detector proposals improve object
enumeration and linkage. Agent Review supplies only a small additional
model-based visual gain at substantially higher cost, and that gain saturates
after one round. These results support a practical two-tier design: use T2 as
the default structured writer and route only uncertain or high-risk cases to a
single review round.

\section*{Data Availability}

The 8,097-image agricultural dataset is closed and is not publicly released.
Raw or processed images, annotations, detection and VLM splits, Golden Test
inputs and references, and per-sample derived artifacts are not distributed
through GitHub or Hugging Face. Aggregate metrics and report materials are
available with the code repository. Access to the public code, report, and
trained checkpoints does not grant access to or permission to redistribute the
underlying dataset.

\section*{Code Availability}

Training, inference, evaluation, audit, paired-bootstrap, and release-building
code is publicly available at
\url{https://github.com/Leo-Yuyang/cv801_final_project}. Project-trained
checkpoints and checkpoint-level metadata are available at
\url{https://huggingface.co/leo20000306/cv801_models}. Public upstream
checkpoints are not mirrored and must be obtained from their original
publishers under the applicable licence and access terms.

\bibliography{references}

\end{document}
